\documentclass[10pt]{article}
\pagestyle{empty}

% --- tight margins (no geometry dependency)
\setlength{\topmargin}{-0.6in}
\setlength{\textheight}{10.3in}
\setlength{\oddsidemargin}{-0.25in}
\setlength{\evensidemargin}{-0.25in}
\setlength{\textwidth}{7.0in}
\setlength{\parindent}{0pt}
\setlength{\parskip}{6pt}

\begin{document}

{\LARGE \textbf{Aaron Beavers}}\\
West Des Moines, IA 50265 \;\;\textbar\;\; 515-388-0539 \;\;\textbar\;\; foamy125@gmail.com\\
GitHub: github.com/riatzukiza

\vspace{10pt}
To the OpenAI team,

Parameter Golf is exactly the kind of problem I want to spend time on: hard constraints, a clean metric, and a lot of room for creative engineering. The part that hooks me is not just ``train a tiny model,'' but the full shape of the challenge --- architecture, compression, reproducibility, measurement discipline, and the ability to prove that an improvement is real.

My strongest evidence is not that I have already trained frontier-scale models. It is that I repeatedly build the systems around hard model problems that make serious iteration possible:
\begin{itemize}
  \item \textbf{Shibboleth}: a generative Clojure DSL and 7-stage evaluation pipeline for adversarial prompt datasets, with leakage-proof splits, provenance, manifests, and reproducibility bundles.
  \item \textbf{Open Hax Proxy}: OpenAI-compatible model infrastructure spanning multiple providers, with routing, reasoning-trace preservation, and tooling that makes experiments easier to run and inspect.
  \item \textbf{Gates of Aker}: simulation and observability work with instrumented systems, clear control surfaces, and tests for complex behavior over time.
\end{itemize}

What I would bring to Parameter Golf is a builder/research-engineer mindset:
\begin{itemize}
  \item disciplined experiment loops instead of vibes,
  \item reproducible logs and artifacts instead of hand-wavy wins,
  \item and a willingness to chase weird ideas so long as they can be measured honestly.
\end{itemize}

I am already joining the challenge. If my work stands out, I would be glad to talk about contributing more directly at OpenAI.

--- Aaron Beavers

\end{document}
